% ---------------------------------------------------------------------------
% Chapter 6: The Children Produce Grandchildren: B-Hadron Decays
% Generated from the outline; edit freely, this file is now the source.
% ---------------------------------------------------------------------------
\chapter{The Children Produce Grandchildren: B-Hadron Decays}
\chaptermark{The Children Produce Grandchildren}
\label{ch:children_producing_grandchildren}

\begin{journeybox}
The B~hadron lives long enough to walk a few millimetres on its own before it
decays, and that walk is the family's birthmark. Then come the grandchildren:
a charm hadron, a few pions, and sometimes a muon with a neutrino. The muon
walks ahead of everyone and the neutrino emigrates without a trace, so for the
first time on the journey energy leaves the family for good.
\end{journeybox}

\physicsline{B-hadron lifetimes and flight distance, displaced vertices, B-hadron decay tables, EvtGen, semileptonic decays, neutrino and soft-muon losses}

\section{B-hadron lifetimes and displaced vertices: the family's birthmark}
\label{sec:children_producing_grandchildren:b_hadron_lifetimes_and_displaced_vertice}
\todo{Write this section.}

\section{Semileptonic decays: neutrinos and soft muons leave the family}
\label{sec:children_producing_grandchildren:semileptonic_decays_neutrinos_and_soft_m}
\todo{Write this section.}

\section{Our jet at this stage: GenJet with and without neutrinos}
\label{sec:children_producing_grandchildren:our_jet_at_this_stage_genjet_with_and_wi}
\todo{Numbers, plots and event display of the $b$-jet at this stage.}

\begin{ledger}{The Children Produce Grandchildren}
  \ledgerrow{before this stage}{}{}{}{--}
  \ledgerrow{after this stage}{}{}{}{}
\end{ledger}

\section{Further reading}
\label{sec:children_producing_grandchildren:further_reading}
Official and theory references for this stage: \cite{Lange:2001uf}, \cite{ALEPH:2001pfo}.

\begin{pitfalls}
\todo{Pitfalls and FAQs for newcomers at this stage.}
\end{pitfalls}

\begin{exercises}
\item \todo{Exercise 1.}
\item \todo{Exercise 2.}
\end{exercises}
